% T19 source: discussion.tex lines 1--56
The results support interpreting the model not as a formula in which ``the
number of agents divides the completion time,'' but as a planning model for a
capacity-constrained human--agent cell.

\subsection{Bound, Optimum, and a Concrete Schedule}

$B_4$ identifies a time below which completion is impossible because of the
aggregate workload of the streams, the critical path induced by the original
precedence constraints, or the single human resource. $T^*$ answers a
different question: the best makespan after accounting for all resource
orderings. Finally, $T(\pi)$ characterizes a selected policy or an actual
schedule. The equality $T^*=B_4$ requires a certificate; otherwise, an equal
lower bound does not imply an equal makespan.

\subsection{Resource Chains as a Gap Mechanism}

Repeated requests from multiple tasks to a single developer link phases that
do not lie on the same path in the original task-DAG\@. Retention of an agent slot
while waiting propagates this competition further through the schedule. The
resource-augmented phase graph makes the mechanism visible: its maximum-weight
path contains not only technological arcs but also resource-order arcs. This
explains why $A,H,W_4,L_4,B_4$ alone do not identify $T^*$.

\subsection{Parallelism as a Configuration Choice}

Available capacity must be distinguished from configured parallelism. With
durations held constant, an additional stream may be left idle, so the optimum
cannot worsen. An organizational transition to a larger $P$, however, may
increase the integration cost $C(P)$ and the switching cost $\gamma(P)$. The
comparison then concerns different processes, and a finite optimum is possible.
In practical terms, $P$ cannot be selected independently of the mode assignment
and the overhead functions. The formal configuration $\sigma=(P,\rho)$ fixes
the number of streams and the task modes; decomposition, the DAG, the phase
profiles, and the acceptance criterion constitute the remaining conditions of
a fully specified scenario. If at least one of them changes together with $P$,
the observed difference pertains to a comparison of the scenarios as a whole,
not to the isolated effect of an additional stream.

\subsection{Decomposition and Parallelism}

Decomposition is beneficial while it shortens long sequential segments and
creates ready work for available streams. Beyond that point, additional
boundaries increase task-specification and integration overhead and the number
of human interventions. Under the decomposition operator studied here, the
optimal granularity therefore depends on $P$: without exploitable parallelism,
small independent tasks provide little benefit, whereas at large $P$,
granularity itself may constrain capacity utilization. The experiments support
this complementarity on the selected grids but do not establish it for an
arbitrary DAG or splitting rule; its direction and magnitude require
calibration on real processes.

\subsection{Volume and Placement of Human Work}

Total $H$ imposes an unconditional sequential workload and a ceiling on
speedup. Its placement is a separate characteristic. When agent and
human-attention phases are scaled equally, moving $H$ among tasks need not
change the aggregate lower bound; when $\gamma>C$, human work on the critical
path is more costly than agent work and lengthens that path. Even when $L_4$
remains unchanged, phase order may change the queue and $T^*$. Thus, process
design requires more than reducing the hours spent on review: it also requires
deciding when and in which tasks those hours are needed.
